\chapter{Sheaves and Sheaf Cohomology}\label{sheaves}
\chapterstatus{complete}

\section{Introduction}\label{sheaves:section-introduction}

A sheaf on a space records data that are defined locally and glue uniquely: continuous, smooth or holomorphic functions,
sections of a bundle, solutions of a differential equation. Sheaves are the language in which complex and algebraic geometry and
Hodge theory are written: the cohomology of sheaves (Section~\ref{sheaves:section-sheaf-cohomology}) unifies singular, de Rham and Dolbeault
cohomology, and the comparison between them is the starting point of Hodge theory. This chapter develops the basic theory: sheaves
and stalks, sheafification, the abelian category of sheaves, direct and inverse images, local systems, and sheaves of modules.
References are \cite{godement-faisceaux}, \cite{bredon-sheaf}, \cite{ks-sheaves-manifolds}, \cite[Chapter II]{hartshorne-ag} and
\cite{stacks-project}.

\noindent\textbf{Prerequisites.} Chapter~\ref{general-topology} (General Topology), Chapter~\ref{categories} (Categories),
Chapter~\ref{homological-algebra} (Homological Algebra) and Chapter~\ref{homotopy} (Homotopy and Covering Spaces).

Throughout, $X$ is a topological space and $\mathrm{Open}(X)$ is the set of open subsets of $X$, ordered by inclusion and viewed as a category.

\section{Presheaves and sheaves}\label{sheaves:section-sheaves}

\begin{definition}\label{sheaves:definition-presheaf}
A \emph{presheaf} $\mathcal{F}$ on $X$ with values in a category $\mathcal{C}$ (sets, abelian groups, rings, modules) is a functor
$\mathcal{F} \colon \mathrm{Open}(X)^{\op} \to \mathcal{C}$: an object $\mathcal{F}(U)$, the \emph{sections} over $U$, for each open $U$, and \emph{restriction maps}
$\mathcal{F}(U) \to \mathcal{F}(V)$, $s \mapsto s|_V$, for $V \subseteq U$, with $s|_U = s$ and $(s|_V)|_W = s|_W$. We also write $\Gamma(U, \mathcal{F}) = \mathcal{F}(U)$. A
\emph{morphism} of presheaves is a natural transformation. A presheaf of sets (or abelian groups, \ldots) is a \emph{sheaf} if for every open cover
$U = \bigcup_i U_i$:
\begin{enumerate}
\item (Locality) if $s, t \in \mathcal{F}(U)$ satisfy $s|_{U_i} = t|_{U_i}$ for all $i$, then $s = t$;
\item (Gluing) if $s_i \in \mathcal{F}(U_i)$ satisfy $s_i|_{U_i \cap U_j} = s_j|_{U_i \cap U_j}$ for all $i, j$, there is $s \in \mathcal{F}(U)$ with $s|_{U_i} = s_i$ for all $i$.
\end{enumerate}
Sheaves form full subcategories $\mathrm{Sh}(X) \subseteq \mathrm{PSh}(X)$ and $\mathrm{Ab}(X) \subseteq \mathrm{PAb}(X)$. Applied to the empty cover of $\emptyset$, the axioms say that
$\mathcal{F}(\emptyset)$ is a single point (the zero group).
\end{definition}

\begin{example}\label{sheaves:example-sheaves}
\begin{enumerate}
\item For a space $Y$, $U \mapsto \{\text{continuous maps } U \to Y\}$ is a sheaf of sets, because continuity is local (pasting lemma,
Lemma~\ref{general-topology:lemma-pasting}). For $Y = \RR$ or $\CC$ it is a sheaf of rings $\mathcal{C}_X$. Likewise smooth functions on a manifold and holomorphic
functions on an open subset of $\CC$ form sheaves of rings.
\item For a continuous map $p \colon E \to X$, the continuous sections $s \colon U \to E$ with $ps = \id_U$ form a sheaf of sets, the \emph{sheaf of sections} of $p$.
\item Let $A$ be a set with at least two elements. The \emph{constant presheaf} $U \mapsto A$ (with $\emptyset \mapsto$ a point) is not a sheaf if $X$ has two disjoint nonempty open
subsets $U_1, U_2$: sections $a \neq b$ on them do not glue to a section on $U_1 \sqcup U_2$. The \emph{constant sheaf} $A_X$ is $U \mapsto \{\text{locally constant maps } U \to A\}$.
\item For $x \in X$ and a set (group) $A$, the \emph{skyscraper sheaf} $x_*A$ has $x_*A(U) = A$ if $x \in U$ and a point otherwise.
\item The presheaf of bounded continuous functions on $\RR$ satisfies locality but not gluing.
\end{enumerate}
\end{example}

\begin{definition}\label{sheaves:definition-stalk}
The \emph{stalk} of a presheaf $\mathcal{F}$ at $x$ is $\mathcal{F}_x = \colim_{U \ni x}\mathcal{F}(U)$, the colimit over the open neighbourhoods of $x$ directed by reverse
inclusion (Example~\ref{categories:example-directed}). By Proposition~\ref{categories:proposition-filtered-colimits-sets}, its elements are \emph{germs} $s_x$ of sections
$s \in \mathcal{F}(U)$, $x \in U$, where $s_x = t_x$ for $t \in \mathcal{F}(V)$ if and only if $s|_W = t|_W$ for some open $W \ni x$ in $U \cap V$. For presheaves of abelian groups
the stalk is an abelian group (Lemma~\ref{poincare-duality:lemma-directed-colimits}). A morphism $\varphi$ induces maps $\varphi_x \colon \mathcal{F}_x \to \mathcal{G}_x$.
\end{definition}

\begin{proposition}\label{sheaves:proposition-stalks-detect}
Let $\varphi \colon \mathcal{F} \to \mathcal{G}$ be a morphism of sheaves.
\begin{enumerate}
\item Two sections $s, t \in \mathcal{F}(U)$ with $s_x = t_x$ for all $x \in U$ are equal.
\item $\varphi_U$ is injective for all $U$ if and only if $\varphi_x$ is injective for all $x$.
\item $\varphi$ is an isomorphism if and only if $\varphi_x$ is bijective for all $x$.
\end{enumerate}
\end{proposition}

\begin{proof}
(1) Each $x \in U$ has an open neighbourhood $W_x \subseteq U$ with $s|_{W_x} = t|_{W_x}$; apply locality to the cover $(W_x)$. (2) If all $\varphi_x$ are injective and
$\varphi(s) = \varphi(t)$, then $\varphi(s)_x = \varphi(t)_x$ gives $s_x = t_x$ for all $x$, so $s = t$ by (1). Conversely, if $\varphi_x(s_x) = \varphi_x(t_x)$ with $s, t \in \mathcal{F}(U)$,
then $\varphi(s)|_W = \varphi(t)|_W$ for some $W \ni x$, so $s|_W = t|_W$ and $s_x = t_x$. (3) An isomorphism induces bijections on stalks. Conversely, by (2) each $\varphi_U$ is
injective. Let $t \in \mathcal{G}(U)$. For each $x \in U$ choose $s^{(x)} \in \mathcal{F}(W_x)$ with $\varphi(s^{(x)})_x = t_x$; shrinking $W_x$, $\varphi(s^{(x)}) = t|_{W_x}$. On
$W_x \cap W_y$, $\varphi(s^{(x)}) = \varphi(s^{(y)})$, so $s^{(x)} = s^{(y)}$ there by injectivity. By gluing there is $s$ with $s|_{W_x} = s^{(x)}$, and $\varphi(s) = t$ by locality.
\end{proof}

\begin{proposition}[Gluing sheaves]\label{sheaves:proposition-gluing}
Let $(U_i)_{i \in I}$ be an open cover of $X$, let $\mathcal{F}_i$ be a sheaf on $U_i$, and write $U_{ij} = U_i \cap U_j$, $U_{ijk} = U_i \cap U_j \cap U_k$. Let
$\varphi_{ij} \colon \mathcal{F}_j|_{U_{ij}} \to \mathcal{F}_i|_{U_{ij}}$ be isomorphisms with $\varphi_{ii} = \id$ and $\varphi_{ij} \circ \varphi_{jk} = \varphi_{ik}$ on $U_{ijk}$ (the \emph{cocycle condition}). Then there is a sheaf $\mathcal{F}$
on $X$ with isomorphisms $\psi_i \colon \mathcal{F}|_{U_i} \to \mathcal{F}_i$ such that $\psi_i = \varphi_{ij} \circ \psi_j$ on $U_{ij}$, and $(\mathcal{F}, (\psi_i))$ is unique up to unique isomorphism.
\end{proposition}

\begin{proof}
For $U$ open let $\mathcal{F}(U)$ be the set of families $(s_i)_i$ with $s_i \in \mathcal{F}_i(U \cap U_i)$ and $s_i = \varphi_{ij}(s_j)$ on $U \cap U_{ij}$ for all $i, j$, with componentwise restriction. The sheaf
axioms hold componentwise, since each $\mathcal{F}_i$ is a sheaf and the compatibility conditions can be checked locally (by locality of the $\mathcal{F}_i$). For $V \subseteq U_i$ define
$\psi_i(( s_j)_j) = s_i$. Its inverse sends $t \in \mathcal{F}_i(V)$ to the family $s_j = \varphi_{ji}(t|_{V \cap U_j})$; this family is compatible because
$\varphi_{jk}(s_k) = \varphi_{jk}\varphi_{ki}(t) = \varphi_{ji}(t) = s_j$, and $s_i = t$ because $\varphi_{ii} = \id$. Conversely a compatible family is determined by its $i$-th member, as $s_j = \varphi_{ji}(s_i)$; so
$\psi_i$ is an isomorphism, and $\psi_i = \varphi_{ij}\psi_j$ is the compatibility condition. If $(\mathcal{G}, (\psi'_i))$ is another solution, the isomorphisms $\psi_i^{-1}\psi'_i \colon \mathcal{G}|_{U_i} \to \mathcal{F}|_{U_i}$ agree on
the overlaps, since $\psi_i^{-1}\psi'_i = \psi_j^{-1}\varphi_{ij}^{-1}\varphi_{ij}\psi'_j = \psi_j^{-1}\psi'_j$; morphisms of sheaves that agree on overlaps glue uniquely, by the sheaf axioms of the target applied
to the images of sections, so they define the unique isomorphism $\mathcal{G} \to \mathcal{F}$ compatible with the $\psi$'s.
\end{proof}

\begin{example}\label{sheaves:example-gluing}
Vector bundles are the basic instance: on a cover trivialising a bundle, the sheaves of sections $\mathcal{F}_i \cong \mathcal{C}^r_{U_i}$ are glued by the transition functions, and the cocycle condition of
Proposition~\ref{sheaves:proposition-gluing} is the cocycle condition for transition functions. Local systems (Section~\ref{sheaves:section-local-systems}) are glued in the same way from constant sheaves, with locally
constant transition functions.
\end{example}

\section{Sheafification}\label{sheaves:section-sheafification}

\begin{theorem}\label{sheaves:theorem-sheafification}
Let $\mathcal{F}$ be a presheaf of sets (abelian groups). For $U$ open let $\mathcal{F}^\#(U)$ be the set of functions $s$ assigning to each $x \in U$ a germ $s(x) \in \mathcal{F}_x$
such that every $x \in U$ has an open neighbourhood $W \subseteq U$ and $t \in \mathcal{F}(W)$ with $s(y) = t_y$ for all $y \in W$. With restriction of functions,
$\mathcal{F}^\#$ is a sheaf, the \emph{sheafification} of $\mathcal{F}$, and $\theta \colon \mathcal{F} \to \mathcal{F}^\#$, $t \mapsto (y \mapsto t_y)$, has the following properties.
\begin{enumerate}
\item $\theta_x \colon \mathcal{F}_x \to \mathcal{F}^\#_x$ is bijective for all $x$.
\item Every morphism $\varphi \colon \mathcal{F} \to \mathcal{G}$ to a sheaf factors uniquely as $\varphi = \psi \circ \theta$. So sheafification is left adjoint to the inclusion of sheaves
into presheaves.
\item If $\mathcal{F}$ is a sheaf, $\theta$ is an isomorphism.
\end{enumerate}
\end{theorem}

\begin{proof}
The conditions defining $\mathcal{F}^\#(U)$ are local, so $\mathcal{F}^\#$ is a sheaf; for abelian groups, the pointwise operations preserve the condition.

(1) Surjectivity: an element of $\mathcal{F}^\#_x$ is the germ of some $s \in \mathcal{F}^\#(U)$, which near $x$ equals $\theta(t)$, so the germ is $\theta_x(t_x)$. Injectivity: if
$\theta(t)_x = \theta(t')_x$, then in particular $t_x = \theta(t)(x) = \theta(t')(x) = t'_x$.

(2) Let $s \in \mathcal{F}^\#(U)$, and choose an open cover $(W_i)$ of $U$ and $t_i \in \mathcal{F}(W_i)$ with $s|_{W_i} = \theta(t_i)$. On $W_i \cap W_j$ the germs of $t_i$ and $t_j$
agree, hence so do the germs of $\varphi(t_i)$ and $\varphi(t_j)$, and by Proposition~\ref{sheaves:proposition-stalks-detect}(1) $\varphi(t_i) = \varphi(t_j)$ there. Let $\psi(s)$ be the
section of $\mathcal{G}$ obtained by gluing. It does not depend on the choices (two choices agree on a common refinement), and $\psi$ is a morphism with $\psi\theta = \varphi$.
Uniqueness: any $\psi'$ with $\psi'\theta = \varphi$ agrees with $\psi$ on each $\theta(t_i)$, hence on $s$ by locality.

(3) By (1) and Proposition~\ref{sheaves:proposition-stalks-detect}(3), applied to the morphism $\theta$ of sheaves.
\end{proof}

\begin{proposition}\label{sheaves:proposition-abelian-category}
$\mathrm{Ab}(X)$ is an abelian category. The kernel of $\varphi \colon \mathcal{F} \to \mathcal{G}$ is $U \mapsto \ker\varphi_U$; its cokernel is the sheafification of the presheaf
$U \mapsto \operatorname{coker}\varphi_U$. Taking stalks is exact: $(\ker\varphi)_x = \ker\varphi_x$ and $(\operatorname{coker}\varphi)_x = \operatorname{coker}\varphi_x$. A sequence of sheaves is exact if
and only if it is exact on all stalks. The functor $\Gamma(U, -)$ is left exact but not exact in general.
\end{proposition}

\begin{proof}
Sums of morphisms and finite direct sums $(\mathcal{F} \oplus \mathcal{G})(U) = \mathcal{F}(U) \oplus \mathcal{G}(U)$ make $\mathrm{Ab}(X)$ additive. The presheaf kernel is a sheaf and is
a kernel in $\mathrm{Ab}(X)$. By Theorem~\ref{sheaves:theorem-sheafification}(2), the sheafified presheaf cokernel is a cokernel. Stalks are directed colimits, which are exact
(Lemma~\ref{poincare-duality:lemma-directed-colimits}), and sheafification does not change stalks; this gives the formulas for stalks. If $\varphi$ is a monomorphism, the
canonical morphism from $\mathcal{F}$ to the kernel of its cokernel is bijective on stalks, because $\mathrm{Ab}$ is abelian; so it is an isomorphism by
Proposition~\ref{sheaves:proposition-stalks-detect}; similarly for epimorphisms. Images and exactness are defined by kernels and cokernels, so they are computed on stalks. Left
exactness of $\Gamma$ holds because kernels are computed sectionwise.

For the failure of right exactness, let $X = \RR$, $p \neq q$ in $X$, and consider $\ZZ_X \to p_*\ZZ \oplus q_*\ZZ$, sending a locally constant function to its values at $p$ and $q$ (when
defined). On stalks it is surjective: at $p$ it is $\ZZ \to \ZZ$, the identity, and at points $x \neq p, q$ the target stalk is $0$, since $x$ has a neighbourhood containing neither
$p$ nor $q$. On global sections it is $\ZZ \to \ZZ \oplus \ZZ$, $n \mapsto (n, n)$, because $\RR$ is connected, and this is not surjective.
\end{proof}

\begin{counterexample}\label{sheaves:counterexample-lifting}
Surjectivity on stalks is a local lifting property, and global sections can fail to lift for topological reasons. Let $T = \RR/\ZZ$ with quotient map $q \colon \RR \to T$. The covering
$p \colon \RR \to S^1$ of Proposition~\ref{homotopy:proposition-circle-covering} induces a continuous bijection $\RR/4\ZZ \to S^1$ from a compact space to a Hausdorff space, hence a homeomorphism
(Proposition~\ref{general-topology:proposition-compact-basic}), so $T \cong \RR/4\ZZ \cong S^1$ and $q$ is a covering map. On any space $X$, let $\mathcal{C}_\RR$ and $\mathcal{C}_T$ be the sheaves of continuous maps
to $\RR$ and to the abelian group $T$, and consider
\[
0 \to \ZZ_X \to \mathcal{C}_\RR \xrightarrow{\ f \mapsto q \circ f\ } \mathcal{C}_T \to 0 .
\]
The kernel of $f \mapsto q \circ f$ consists of the continuous maps to $\ZZ$, that is, $\ZZ_X$. On stalks the map is surjective: a continuous $g \colon W \to T$ near $x$ takes a small neighbourhood of $x$ into an evenly covered
open set, over which $q$ has a continuous inverse on each sheet. So the sequence is exact (Proposition~\ref{sheaves:proposition-abelian-category}). But for $X = T$, the identity $\id_T \in \mathcal{C}_T(T)$ is not
of the form $q \circ f$: a continuous $f \colon T \to \RR$ with $q \circ f = \id$ would be a lift of $\id_T$ through $q$, which by the lifting criterion (Proposition~\ref{homotopy:proposition-lifting-criterion})
requires $\pi_1(T) \subseteq q_*\pi_1(\RR) = 0$, whereas $\pi_1(T) \cong \pi_1(S^1) \cong \ZZ$ (Theorem~\ref{homotopy:theorem-circle}). So $\Gamma(T, -)$ is not right exact on this sequence; the obstruction
will reappear as the nonzero group $H^1(S^1, \ZZ) \cong \ZZ$ (Section~\ref{sheaves:section-sheaf-cohomology}).
\end{counterexample}

\section{Direct and inverse images}\label{sheaves:section-images}

\begin{definition}\label{sheaves:definition-images}
Let $f \colon X \to Y$ be continuous. The \emph{direct image} of a sheaf $\mathcal{F}$ on $X$ is the sheaf $f_*\mathcal{F}(V) = \mathcal{F}(f^{-1}V)$ on $Y$. The \emph{inverse image} of a sheaf
$\mathcal{G}$ on $Y$ is the sheafification $f^{-1}\mathcal{G}$ of the presheaf $f^{-1}_p\mathcal{G}(U) = \colim_{V \supseteq f(U)}\mathcal{G}(V)$, the colimit over open $V \supseteq f(U)$
directed by reverse inclusion. For $i \colon Z \to X$ the inclusion of a subspace, $\mathcal{F}|_Z = i^{-1}\mathcal{F}$ is the \emph{restriction}; for $Z$ open,
$\mathcal{F}|_Z(W) = \mathcal{F}(W)$.
\end{definition}

\begin{proposition}\label{sheaves:proposition-images}
\begin{enumerate}
\item $(f^{-1}\mathcal{G})_x \cong \mathcal{G}_{f(x)}$; hence $f^{-1}$ is exact on sheaves of abelian groups. The functor $f_*$ is left exact.
\item There is a bijection $\Hom_{\mathrm{Sh}(X)}(f^{-1}\mathcal{G}, \mathcal{F}) \cong \Hom_{\mathrm{Sh}(Y)}(\mathcal{G}, f_*\mathcal{F})$, natural in $\mathcal{F}$ and $\mathcal{G}$: $f^{-1}$ is left adjoint to
$f_*$.
\item $(g \circ f)_* = g_* \circ f_*$ and $(g \circ f)^{-1} \cong f^{-1} \circ g^{-1}$.
\end{enumerate}
\end{proposition}

\begin{proof}
(1) The pairs $(U, V)$ with $x \in U$ and $V \supseteq f(U)$ map cofinally to the open neighbourhoods $V$ of $f(x)$ (take $U = f^{-1}V$), so
$(f^{-1}_p\mathcal{G})_x = \colim_{V \ni f(x)}\mathcal{G}(V)$ by Lemma~\ref{poincare-duality:lemma-directed-colimits}, and sheafification preserves stalks. Exactness follows from
Proposition~\ref{sheaves:proposition-abelian-category}. Kernels of sheaves are computed sectionwise, so $f_*$ preserves them.

(2) By Theorem~\ref{sheaves:theorem-sheafification}(2), morphisms $f^{-1}\mathcal{G} \to \mathcal{F}$ are morphisms of presheaves $f_p^{-1}\mathcal{G} \to \mathcal{F}$, that is, families of
maps $\mathcal{G}(V) \to \mathcal{F}(U)$ for $V \supseteq f(U)$, compatible with restrictions in $U$ and $V$. Such a family determines maps $\mathcal{G}(V) \to \mathcal{F}(f^{-1}V)$ (take
$U = f^{-1}V$), compatible with restrictions, that is, a morphism $\mathcal{G} \to f_*\mathcal{F}$. Conversely a morphism $\psi \colon \mathcal{G} \to f_*\mathcal{F}$ gives
$\mathcal{G}(V) \xrightarrow{\psi_V} \mathcal{F}(f^{-1}V) \to \mathcal{F}(U)$ by restriction, as $U \subseteq f^{-1}V$. These constructions are mutually inverse and natural.

(3) The first equality is clear; the second follows from (2) and uniqueness of adjoints (Proposition~\ref{categories:proposition-adjoint-unique}).
\end{proof}

\begin{definition}\label{sheaves:definition-extension-by-zero}
Let $j \colon U \to X$ be the inclusion of an open subset and $\mathcal{F}$ a sheaf of abelian groups on $U$. The \emph{extension by zero} $j_!\mathcal{F}$ is the sheafification of the
presheaf $W \mapsto \mathcal{F}(W)$ if $W \subseteq U$ and $W \mapsto 0$ otherwise.
\end{definition}

\begin{proposition}\label{sheaves:proposition-extension-by-zero}
$(j_!\mathcal{F})_x = \mathcal{F}_x$ for $x \in U$ and $0$ for $x \notin U$. If $Z = X \setminus U$ with inclusion $i$, then for every sheaf of abelian groups $\mathcal{F}$ on $X$ there is a
short exact sequence $0 \to j_!(\mathcal{F}|_U) \to \mathcal{F} \to i_*(\mathcal{F}|_Z) \to 0$.
\end{proposition}

\begin{proof}
For $x \in U$ the open neighbourhoods contained in $U$ are cofinal; for $x \notin U$ no neighbourhood is contained in $U$. The first map is induced, via
Theorem~\ref{sheaves:theorem-sheafification}(2), by the inclusions $\mathcal{F}(W) \to \mathcal{F}(W)$ for $W \subseteq U$; the second is the unit
$\mathcal{F} \to i_*i^{-1}\mathcal{F}$ of the adjunction. On stalks at $x \in U$, the sequence is $0 \to \mathcal{F}_x \to \mathcal{F}_x \to 0 \to 0$, since $x$ has a neighbourhood missing the
closed set $Z$. At $z \in Z$, the open subsets $W \cap Z$ are cofinal among neighbourhoods of $z$ in $Z$, so $(i_*i^{-1}\mathcal{F})_z = (i^{-1}\mathcal{F})_z = \mathcal{F}_z$, and the
sequence is $0 \to 0 \to \mathcal{F}_z \to \mathcal{F}_z \to 0$.
\end{proof}

\section{Locally constant sheaves and local systems}\label{sheaves:section-local-systems}

\begin{definition}\label{sheaves:definition-local-system}
A sheaf $\mathcal{F}$ of sets is \emph{locally constant} if every point has an open neighbourhood $U$ with $\mathcal{F}|_U$ isomorphic to a constant sheaf $A_U$. A \emph{local system}
of $K$-vector spaces (for a field $K$) is a locally constant sheaf of $K$-vector spaces whose stalks are finite-dimensional; local systems of abelian groups or modules are defined
similarly.
\end{definition}

\begin{proposition}\label{sheaves:proposition-etale-space}
Let $\mathcal{F}$ be a sheaf of sets on $X$. The \emph{\'etal\'e space} $E_\mathcal{F} = \bigsqcup_x\mathcal{F}_x$, with the topology whose basis is the sets $\{s_x : x \in U\}$ for
$s \in \mathcal{F}(U)$, has a local homeomorphism $p \colon E_\mathcal{F} \to X$, and $\mathcal{F}$ is isomorphic to the sheaf of sections of $p$. If $\mathcal{F}$ is locally constant and $X$
is locally connected, $p$ is a covering map (except that fibres may be empty), and conversely the sheaf of sections of such a covering is locally constant.
\end{proposition}

\begin{proof}
If $s_x = t_x$ for $s \in \mathcal{F}(U)$, $t \in \mathcal{F}(V)$, the two basic sets contain the basic set of $s|_W$ for some $W \ni x$; so these sets form a basis, and $p$ maps the basic
set of $s$ homeomorphically onto $U$, with inverse $\hat{s} \colon x \mapsto s_x$. The map $s \mapsto \hat{s}$ from $\mathcal{F}(U)$ to continuous sections over $U$ is injective by
Proposition~\ref{sheaves:proposition-stalks-detect}(1); it is surjective, because a continuous section $\sigma$ is locally of the form $\hat{s}$ (near $x$, $\sigma$ lies in the basic
set of some $s$ with $s_x = \sigma(x)$), and these local $s$ glue. If $\mathcal{F}|_U \cong A_U$ with $U$ connected, then $p^{-1}(U) \cong U \times A$ with $A$ discrete, because the germs of
the constant sections $a \in A$ give disjoint open sets covering $p^{-1}(U)$. Conversely, over a connected evenly covered open set, sections of a covering correspond to locally constant
choices of sheets.
\end{proof}

\begin{theorem}\label{sheaves:theorem-local-systems-representations}
Let $X$ be path-connected, locally path-connected and semilocally simply connected, and $x \in X$. Taking the stalk at $x$, with the monodromy action, gives equivalences of
categories between
\begin{enumerate}
\item locally constant sheaves of sets on $X$ and sets with a right action of $\pi_1(X, x)$;
\item local systems of $K$-vector spaces on $X$ and finite-dimensional right $K[\pi_1(X, x)]$-modules, that is, representations
$\rho \colon \pi_1(X, x)^{\op} \to \GL(\mathcal{F}_x)$.
\end{enumerate}
\end{theorem}

\begin{proof}[Sketch of proof]
By Proposition~\ref{sheaves:proposition-etale-space}, locally constant sheaves correspond to covering spaces (with possibly empty fibres), and morphisms of sheaves to maps of
coverings. A covering is the disjoint union of its components, which are connected coverings of $X$, and the fibre over $x$ with its monodromy action
(Corollary~\ref{homotopy:corollary-lifting-consequences}) is a right $\pi_1(X, x)$-set whose orbits correspond to the components. For connected coverings, the classification
theorem (Theorem~\ref{homotopy:theorem-classification}) identifies isomorphism classes with transitive $\pi_1$-sets $H\backslash\pi_1(X, x)$, and the lifting criterion
(Proposition~\ref{homotopy:proposition-lifting-criterion}) identifies maps of coverings with equivariant maps of fibres; this gives (1). The omitted verification is that these
identifications are functorial; see \cite[Chapter 2]{szamuely-galois}. For (2), a local system is a locally constant sheaf of sets with a $K$-vector space structure on the stalks
compatible with monodromy, and the monodromy action is then $K$-linear.
\end{proof}

\begin{example}\label{sheaves:example-local-systems}
\begin{enumerate}
\item On a simply connected space every local system is constant.
\item On $S^1$, the rank one local systems of $\CC$-vector spaces correspond to $\lambda \in \CC^\times$, the monodromy around the generator of $\pi_1(S^1) \cong \ZZ$. For
$\lambda \neq 1$ the local system $\mathcal{L}_\lambda$ has no nonzero global sections, since a global section is a vector fixed by monodromy. The real local system with monodromy $-1$ is the
sheaf of locally constant sections of the M\"obius band, viewed as a line bundle over $S^1$.
\item For a manifold $M$, the \emph{orientation sheaf} $\mathit{or}_M$, with stalks $H_n(M \mid x; \ZZ)$, is the local system of rank one associated with the orientation double covering
(Proposition~\ref{poincare-duality:proposition-orientation-covering}); it is constant if and only if $M$ is orientable.
\item For a family of spaces $f \colon E \to B$ that is locally trivial, the cohomology groups of the fibres form local systems $R^kf_*\ZZ$ on $B$ (Section~\ref{sheaves:section-sheaf-cohomology}); their
monodromy is a central object in the theory of variations of Hodge structure (Chapter~\ref{vhs}).
\end{enumerate}
\end{example}

\section{Ringed spaces and sheaves of modules}\label{sheaves:section-ringed-spaces}

\begin{definition}\label{sheaves:definition-ringed-space}
A \emph{ringed space} is a pair $(X, \mathcal{O}_X)$ of a space and a sheaf of commutative rings. It is \emph{locally ringed} if all stalks $\mathcal{O}_{X,x}$ are local rings. A
\emph{morphism} $(f, f^\#) \colon (X, \mathcal{O}_X) \to (Y, \mathcal{O}_Y)$ is a continuous map $f$ with a morphism of sheaves of rings $f^\# \colon \mathcal{O}_Y \to f_*\mathcal{O}_X$; for locally
ringed spaces one requires the induced maps $\mathcal{O}_{Y, f(x)} \to \mathcal{O}_{X, x}$ to be local homomorphisms. An \emph{$\mathcal{O}_X$-module} is a sheaf of abelian groups
$\mathcal{F}$ with $\mathcal{O}_X(U)$-module structures on $\mathcal{F}(U)$ compatible with restrictions.
\end{definition}

\begin{example}\label{sheaves:example-ringed-spaces}
A space with its sheaf of continuous real functions, a smooth manifold with smooth functions (Chapter~\ref{smooth-manifolds}), a complex manifold with holomorphic functions
(Chapter~\ref{complex-manifolds}) and a scheme (Chapter~\ref{schemes}) are locally ringed spaces: in each case the germs vanishing at $x$ form the unique maximal ideal of the stalk, since a
germ not vanishing at $x$ is invertible near $x$.
\end{example}

\begin{definition}\label{sheaves:definition-module-operations}
Let $\mathcal{F}$, $\mathcal{G}$ be $\mathcal{O}_X$-modules. The \emph{tensor product} $\mathcal{F} \otimes_{\mathcal{O}_X} \mathcal{G}$ is the sheafification of
$U \mapsto \mathcal{F}(U) \otimes_{\mathcal{O}_X(U)} \mathcal{G}(U)$. The sheaf $\mathcal{H}om_{\mathcal{O}_X}(\mathcal{F}, \mathcal{G})$ has sections over $U$ the morphisms
$\mathcal{F}|_U \to \mathcal{G}|_U$ of $\mathcal{O}_U$-modules, and $\mathcal{F}^\vee = \mathcal{H}om(\mathcal{F}, \mathcal{O}_X)$ is the \emph{dual}. An $\mathcal{O}_X$-module is
\emph{locally free of rank $r$} if every point has an open neighbourhood $U$ with $\mathcal{F}|_U \cong \mathcal{O}_U^{\oplus r}$, and \emph{invertible} if it is locally free of rank $1$.
For a morphism $f$ of ringed spaces and an $\mathcal{O}_Y$-module $\mathcal{G}$, the \emph{pullback} is $f^*\mathcal{G} = f^{-1}\mathcal{G} \otimes_{f^{-1}\mathcal{O}_Y} \mathcal{O}_X$.
\end{definition}

\begin{proposition}\label{sheaves:proposition-modules}
\begin{enumerate}
\item $\mathcal{H}om(\mathcal{F}, \mathcal{G})$ is a sheaf, and $(\mathcal{F} \otimes \mathcal{G})_x \cong \mathcal{F}_x \otimes_{\mathcal{O}_{X,x}} \mathcal{G}_x$.
\item If $\mathcal{E}$ and $\mathcal{E}'$ are locally free of ranks $r$ and $r'$, then $\mathcal{E} \otimes \mathcal{E}'$ is locally free of rank $rr'$, $\mathcal{E}^\vee$ is locally free of rank
$r$, and the evaluation map $\mathcal{E} \otimes \mathcal{E}^\vee \to \mathcal{O}_X$ is an isomorphism if $r = 1$.
\item The isomorphism classes of invertible sheaves form an abelian group $\Pic(X)$ under $\otimes$, with identity $\mathcal{O}_X$ and inverse $\mathcal{L}^\vee$.
\end{enumerate}
\end{proposition}

\begin{proof}
(1) Morphisms defined on the members of an open cover and agreeing on overlaps glue uniquely, because sections of $\mathcal{G}$ do. For the stalk, sheafification does not change stalks,
so it suffices to compare $\colim_{U \ni x}\mathcal{F}(U) \otimes_{\mathcal{O}_X(U)} \mathcal{G}(U)$ with $\mathcal{F}_x \otimes_{\mathcal{O}_{X,x}} \mathcal{G}_x$. The map
$[s \otimes t] \mapsto s_x \otimes t_x$ is well defined; an inverse sends $s_x \otimes t_x$ to the class of $s|_W \otimes t|_W$ on a common neighbourhood $W$, which is well defined because
every relation defining the tensor product over $\mathcal{O}_{X,x}$ involves finitely many germs and already holds over some neighbourhood of $x$ \cite[Chapter II, Section 5]{hartshorne-ag}.

(2) All statements are local, so we may assume $\mathcal{E} = \mathcal{O}^r$ and $\mathcal{E}' = \mathcal{O}^{r'}$. Then $\mathcal{O}^r \otimes \mathcal{O}^{r'} \cong \mathcal{O}^{rr'}$ (the
presheaf tensor product is already a sheaf in this case), $\mathcal{H}om(\mathcal{O}^r, \mathcal{O}) \cong \mathcal{O}^r$, and for $r = 1$ evaluation $\mathcal{O} \otimes \mathcal{O} \to \mathcal{O}$
is multiplication, an isomorphism. Since the evaluation map is a globally defined morphism, and being an isomorphism can be checked on stalks
(Proposition~\ref{sheaves:proposition-stalks-detect}), it is an isomorphism.

(3) Tensor product is associative and commutative up to natural isomorphism, $\mathcal{O}_X \otimes \mathcal{L} \cong \mathcal{L}$, and (2) gives inverses.
\end{proof}

\begin{remark}\label{sheaves:remark-vector-bundles}
If $E \to X$ is a (real or complex, topological, smooth or holomorphic) vector bundle of rank $r$, its sheaf of sections is a locally free module of rank $r$ over the corresponding
sheaf of functions, and every locally free sheaf arises in this way, uniquely up to isomorphism (Chapter~\ref{vector-bundles}). Under this correspondence, $\Pic(X)$ is the group of line
bundles; for a complex manifold it is computed by sheaf cohomology via the exponential sequence (Chapter~\ref{complex-manifolds}).
\end{remark}

\section{Sheaf Cohomology}\label{sheaves:section-sheaf-cohomology}
The global sections functor $\Gamma(X, -)$ is left exact but not exact (Proposition~\ref{sheaves:proposition-abelian-category}), and
sheaf cohomology $H^i(X, \mathcal{F})$ is its right derived functor. It can be computed by several kinds of acyclic resolutions:
flasque, soft or fine sheaves (Theorem~\ref{sheaves:theorem-soft-acyclic}), \v{C}ech complexes of acyclic covers
(Theorem~\ref{sheaves:theorem-leray}), and hypercohomology of complexes. On reasonable spaces the cohomology of a constant
sheaf is singular cohomology (Theorem~\ref{sheaves:theorem-singular-comparison}); on a manifold the de Rham complex is a fine
resolution of the constant sheaf, which gives the de Rham theorem (Chapter~\ref{de-rham}). References: \cite[Chapter III]{hartshorne-ag},
\cite{godement-faisceaux}, \cite{bredon-sheaf}, \cite{ks-sheaves-manifolds} and \cite[Chapter 4]{voisin-hodge-1}.

\noindent\textbf{Prerequisites.} Chapter~\ref{sheaves} (Sheaves on Topological Spaces), Section~\ref{homological-algebra:section-spectral-sequences} (Spectral Sequences),
Section~\ref{homological-algebra:section-derived-categories} (Derived Categories) and Section~\ref{singular-homology:section-singular-cohomology} (Singular Cohomology).

\section{Cohomology as a derived functor}\label{sheaves:section-derived}

\begin{lemma}\label{sheaves:lemma-adjoints-injectives}
Let $L \colon \mathcal{B} \to \mathcal{A}$ be an exact functor between abelian categories with right adjoint $R \colon \mathcal{A} \to \mathcal{B}$. Then $R$ preserves injective objects.
In particular $f_*$ preserves injective sheaves for every continuous $f$, and so does restriction $j^{-1}$ to an open subset $j \colon U \to X$.
\end{lemma}

\begin{proof}
For $I$ injective, $\Hom_\mathcal{B}(-, R(I)) \cong \Hom_\mathcal{A}(L(-), I)$ is a composite of exact functors, hence exact. For $f_*$ use
Proposition~\ref{sheaves:proposition-images}. For $j^{-1}$ we show that $j_!$ is an exact left adjoint: a morphism $j_!\mathcal{F} \to \mathcal{G}$ is, by
Theorem~\ref{sheaves:theorem-sheafification}, a compatible family of maps $\mathcal{F}(W) \to \mathcal{G}(W)$ for $W \subseteq U$, that is, a morphism
$\mathcal{F} \to \mathcal{G}|_U$; and $j_!$ is exact on stalks (Proposition~\ref{sheaves:proposition-extension-by-zero}).
\end{proof}

\begin{theorem}\label{sheaves:theorem-enough-injectives}
The category $\mathrm{Ab}(X)$ has enough injectives, and so does the category of $\mathcal{O}_X$-modules on a ringed space.
\end{theorem}

\begin{proof}
For each $x$ embed $\mathcal{F}_x$ into an injective abelian group $I_x$ (Theorem~\ref{homological-algebra:theorem-enough-injectives}). The skyscraper sheaf $x_*I_x$ is injective,
since $\Hom(\mathcal{G}, x_*I_x) \cong \Hom(\mathcal{G}_x, I_x)$ is exact in $\mathcal{G}$; this isomorphism is Proposition~\ref{sheaves:proposition-images}(2) and (1) for $i \colon \{x\} \to X$, as $x_*I_x = i_*I_x$ and $i^{-1}\mathcal{G} = \mathcal{G}_x$. A product of injective objects is
injective, and $\mathcal{F} \to \prod_x x_*I_x$ is a monomorphism, being injective on stalks. For $\mathcal{O}_X$-modules use injective $\mathcal{O}_{X,x}$-modules instead.
\end{proof}

\begin{definition}\label{sheaves:definition-sheaf-cohomology}
The \emph{sheaf cohomology} groups are $H^i(X, \mathcal{F}) = R^i\Gamma(X, -)(\mathcal{F})$ (Definition~\ref{homological-algebra:definition-derived-functor}). For a complex
$\mathcal{K}^\bullet$ of sheaves, bounded below, the \emph{hypercohomology} is $\mathbb{H}^i(X, \mathcal{K}^\bullet) = H^i(R\Gamma(X, \mathcal{K}^\bullet))$
(Definition~\ref{homological-algebra:definition-right-derived}). For a continuous $f \colon X \to Y$, the \emph{higher direct images} are $R^if_*\mathcal{F}$.
\end{definition}

\begin{remark}\label{sheaves:remark-cohomology-modules}
For an $\mathcal{O}_X$-module $\mathcal{F}$, derived functors computed in $\mathcal{O}_X$-modules and in $\mathrm{Ab}(X)$ agree, because injective $\mathcal{O}_X$-modules are flasque
(Lemma~\ref{sheaves:lemma-injective-flasque} applies verbatim with $j_!\mathcal{O}_U \to \mathcal{O}_X$) and flasque sheaves are acyclic in both settings.
\end{remark}

\section{Flasque sheaves}\label{sheaves:section-flasque}

\begin{definition}\label{sheaves:definition-flasque}
A sheaf $\mathcal{F}$ is \emph{flasque} if all restriction maps $\mathcal{F}(U) \to \mathcal{F}(V)$ are surjective.
\end{definition}

\begin{lemma}\label{sheaves:lemma-injective-flasque}
Injective sheaves of abelian groups are flasque. Direct images and restrictions to open subsets of flasque sheaves are flasque, and so are products.
\end{lemma}

\begin{proof}
Let $I$ be injective and $V \subseteq U$ open, with inclusions $j_V, j_U$ into $X$. The morphism $j_{V!}\ZZ_V \to j_{U!}\ZZ_U$ is injective on stalks, so
$\Hom(j_{U!}\ZZ_U, I) \to \Hom(j_{V!}\ZZ_V, I)$ is surjective. By the adjunction in the proof of Lemma~\ref{sheaves:lemma-adjoints-injectives},
$\Hom(j_{U!}\ZZ_U, I) = \Hom(\ZZ_U, I|_U) = I(U)$, a morphism from the constant sheaf being determined by the image of $1$; so $I(U) \to I(V)$ is surjective. The other statements are
immediate from the definitions.
\end{proof}

\begin{proposition}\label{sheaves:proposition-flasque}
Let $0 \to \mathcal{F}' \to \mathcal{F} \xrightarrow{\pi} \mathcal{F}'' \to 0$ be exact with $\mathcal{F}'$ flasque. Then $\mathcal{F}(U) \to \mathcal{F}''(U)$ is surjective for all $U$. If also $\mathcal{F}$ is
flasque, so is $\mathcal{F}''$. Flasque sheaves are $\Gamma$-acyclic, so $H^i(X, \mathcal{F})$ can be computed with flasque resolutions.
\end{proposition}

\begin{proof}
We may take $U = X$. Let $s \in \mathcal{F}''(X)$, and consider pairs $(V, t)$ with $V$ open, $t \in \mathcal{F}(V)$ and $\pi(t) = s|_V$, ordered by extension. A chain has an upper bound by gluing,
and it is nonempty, since $s$ lifts locally. Let $(V, t)$ be maximal and suppose $x \notin V$. Choose $W \ni x$ and $t' \in \mathcal{F}(W)$ lifting $s|_W$. On $V \cap W$, $t - t'$ lies in
$\mathcal{F}'(V \cap W)$, which extends to $u \in \mathcal{F}'(X)$ since $\mathcal{F}'$ is flasque. Then $t$ and $t' + u|_W$ agree on $V \cap W$ and glue to a lift over $V \cup W$, contradicting
maximality. If $\mathcal{F}$ is flasque and $s'' \in \mathcal{F}''(V)$, lift it to $\mathcal{F}(V)$, extend to $\mathcal{F}(X)$ and map to $\mathcal{F}''(X)$.

Let $\mathcal{F}$ be flasque and embed it into an injective $I$ with cokernel $\mathcal{G}$, which is flasque by the above and Lemma~\ref{sheaves:lemma-injective-flasque}.
The long exact sequence gives $H^1(X, \mathcal{F}) = \operatorname{coker}(I(X) \to \mathcal{G}(X)) = 0$ and $H^i(X, \mathcal{F}) \cong H^{i-1}(X, \mathcal{G})$ for $i \geq 2$, which vanishes by
induction on $i$ for all flasque sheaves simultaneously. The last statement is Proposition~\ref{homological-algebra:proposition-acyclic-resolutions}.
\end{proof}

\begin{example}\label{sheaves:example-godement}
For every sheaf $\mathcal{F}$, the sheaf $\mathcal{C}^0(\mathcal{F}) = \prod_x x_*\mathcal{F}_x$ of possibly discontinuous sections is flasque, and $\mathcal{F} \to \mathcal{C}^0(\mathcal{F})$ is
injective. Repeating with the cokernel gives the \emph{Godement resolution}, a canonical flasque resolution, functorial and exact in $\mathcal{F}$.
\end{example}

\section{Soft and fine sheaves}\label{sheaves:section-soft}

In this section $X$ is paracompact Hausdorff. For a subset $Z \subseteq X$ we write $\mathcal{F}(Z) = \Gamma(Z, \mathcal{F}|_Z)$; by Proposition~\ref{sheaves:proposition-etale-space}
this is the set of continuous sections over $Z$ of the \'etal\'e space of $\mathcal{F}$. Sections over the members of a locally finite closed cover of $Z$ that agree on overlaps
glue, by the pasting lemma applied locally.

\begin{lemma}\label{sheaves:lemma-closed-sections-extend}
Let $Z \subseteq X$ be closed. Every $s \in \mathcal{F}(Z)$ extends to a section over some open neighbourhood of $Z$.
\end{lemma}

\begin{proof}
Each $z \in Z$ has an open neighbourhood $W_z$ in $X$ and $t_z \in \mathcal{F}(W_z)$ with germs $s(y)$ at all $y \in W_z \cap Z$. Take a locally finite open refinement $(U_i)_{i \in I}$ of the
cover $\{W_z\} \cup \{X \setminus Z\}$, choosing for each $i$ either a section $t_i = t_z|_{U_i}$ (first kind) or $U_i \subseteq X \setminus Z$ (second kind), and closed sets
$\overline{V_i} \subseteq U_i$ with $(V_i)$ still covering (Lemma~\ref{general-topology:lemma-shrinking}, Proposition~\ref{general-topology:proposition-paracompact-normal}). For $x \in X$
let $J(x) = \{i : x \in \overline{V_i}\}$, a finite set. Let $W$ be the set of $x$ such that $J(x)$ contains only indices of the first kind and the germs $(t_i)_x$, $i \in J(x)$, all agree.

$Z \subseteq W$, since for $z \in Z$ all $i \in J(z)$ are of the first kind and $(t_i)_z = s(z)$. $W$ is open: for $x \in W$ choose an open $N \ni x$ that meets only finitely many $\overline{V_i}$,
avoids those with $i \notin J(x)$, lies in $\bigcap_{i \in J(x)}U_i$, and on which the sections $t_i$, $i \in J(x)$, agree; then $J(y) \subseteq J(x)$ for $y \in N$, so $N \subseteq W$. On such an
$N = N_x$, put $t^{(x)} = t_i|_{N_x}$ for any $i \in J(x)$ (nonempty, since the $V_i$ cover). At $y \in N_x \cap N_{x'}$ both $t^{(x)}$ and $t^{(x')}$ have the germ $(t_k)_y$ for $k \in J(y)$, so
they agree, and they glue to $t \in \mathcal{F}(W)$ with $t_z = s(z)$ on $Z$.
\end{proof}

\begin{definition}\label{sheaves:definition-soft-fine}
A sheaf $\mathcal{F}$ on $X$ is \emph{soft} if $\mathcal{F}(X) \to \mathcal{F}(Z)$ is surjective for every closed $Z$. The \emph{support} of an endomorphism $\eta$ of $\mathcal{F}$ is the
closure of $\{x : \eta_x \neq 0\}$. A sheaf of abelian groups is \emph{fine} if for every locally finite open cover $(U_i)$ there are endomorphisms $\eta_i$ with
$\operatorname{supp}\eta_i \subseteq U_i$ and $\sum_i\eta_i = \id$ (a \emph{partition of unity} of $\mathcal{F}$).
\end{definition}

\begin{example}\label{sheaves:example-fine}
The sheaf $\mathcal{C}_X$ of continuous real functions is fine, by partitions of unity (Theorem~\ref{general-topology:theorem-partitions-of-unity}): take $\eta_i$ to be
multiplication by $\psi_i$. More generally every sheaf of modules over $\mathcal{C}_X$ is fine, and on a smooth manifold every module over the sheaf of smooth functions is fine,
using smooth partitions of unity (Chapter~\ref{smooth-manifolds}). The constant sheaf $\ZZ_X$ on a connected space with more than one point is not soft.
\end{example}

\begin{lemma}\label{sheaves:lemma-soft-properties}
\begin{enumerate}
\item Flasque sheaves and fine sheaves are soft.
\item If $0 \to \mathcal{F}' \to \mathcal{F} \xrightarrow{\pi} \mathcal{F}'' \to 0$ is exact and $\mathcal{F}'$ is soft, then $\mathcal{F}(X) \to \mathcal{F}''(X)$ is surjective.
\item If moreover $\mathcal{F}$ is soft, then $\mathcal{F}''$ is soft.
\end{enumerate}
\end{lemma}

\begin{proof}
(1) Let $s \in \mathcal{F}(Z)$ and extend it to an open $W \supseteq Z$ by Lemma~\ref{sheaves:lemma-closed-sections-extend}. If $\mathcal{F}$ is flasque, extend further to $X$. If
$\mathcal{F}$ is fine, take $\eta_1 + \eta_2 = \id$ for the cover $\{W, X \setminus Z\}$. The section $\eta_1(s)$ of $\mathcal{F}(W)$ vanishes on $W \setminus \operatorname{supp}\eta_1$, so it glues
with $0$ on $X \setminus \operatorname{supp}\eta_1$ to a global section $t$; at $z \in Z$, $t_z = s_z - (\eta_2)_z(s_z) = s_z$ since $z \notin \operatorname{supp}\eta_2$.

(2) Choose a locally finite open cover $(U_i)$ with lifts $s_i \in \mathcal{F}(U_i)$ of $s|_{U_i}$ for a given $s \in \mathcal{F}''(X)$, and closed $\overline{V_i} \subseteq U_i$ covering $X$. Consider pairs
$(J, t)$ with $J \subseteq I$ and $t \in \mathcal{F}(Z_J)$ lifting $s$, where $Z_J = \bigcup_{i \in J}\overline{V_i}$, which is closed (Lemma~\ref{general-topology:lemma-locally-finite}). Chains
have upper bounds by gluing along the locally finite closed cover, so Zorn's lemma gives a maximal $(J, t)$. If $i \notin J$, the difference $t - s_i$ on $Z_J \cap \overline{V_i}$ is a section of
$\mathcal{F}'$ (its image under $\pi$ is $0$), which extends to $u \in \mathcal{F}'(X)$. Then $t$ and $s_i + u$ agree on $Z_J \cap \overline{V_i}$ and glue to a lift over $Z_{J \cup \{i\}}$, a
contradiction. So $J = I$.

(3) A closed subspace $Z$ is paracompact Hausdorff, the restricted sequence on $Z$ is exact, and $\mathcal{F}'|_Z$ is soft. By (2) applied on $Z$, $\mathcal{F}(Z) \to \mathcal{F}''(Z)$ is surjective,
and $\mathcal{F}(X) \to \mathcal{F}(Z)$ is surjective.
\end{proof}

\begin{theorem}\label{sheaves:theorem-soft-acyclic}
On a paracompact Hausdorff space, soft sheaves, in particular fine sheaves, are $\Gamma$-acyclic. Hence cohomology can be computed with soft or fine resolutions.
\end{theorem}

\begin{proof}
Identical to the proof of Proposition~\ref{sheaves:proposition-flasque}, using Lemma~\ref{sheaves:lemma-soft-properties} and the fact that injective sheaves are
flasque, hence soft.
\end{proof}

\begin{corollary}\label{sheaves:corollary-fine-acyclic}
Let $X$ be paracompact Hausdorff. Then $H^i(X, \mathcal{F}) = 0$ for $i > 0$ for every sheaf of $\mathcal{C}_X$-modules, for instance the sheaf of sections of a real vector bundle.
\end{corollary}

\begin{proof}
Such a sheaf is fine (Example~\ref{sheaves:example-fine}) and is its own resolution.
\end{proof}

\section{\v{C}ech cohomology}\label{sheaves:section-cech}

\begin{definition}\label{sheaves:definition-cech}
Let $\mathcal{U} = (U_i)_{i \in I}$ be an open cover of $X$ and $\mathcal{F}$ a presheaf of abelian groups. Write $U_{i_0 \cdots i_p} = U_{i_0} \cap \cdots \cap U_{i_p}$. An \emph{alternating
\v{C}ech $p$-cochain} is a family $\alpha = (\alpha_{i_0 \cdots i_p})$, indexed by all $(i_0, \ldots, i_p) \in I^{p+1}$, with $\alpha_{i_0 \cdots i_p} \in \mathcal{F}(U_{i_0 \cdots i_p})$, which
changes sign when two indices are exchanged and vanishes when two indices are equal. They form a group $\check{C}^p(\mathcal{U}, \mathcal{F})$, with differential
\[
(d\alpha)_{i_0 \cdots i_{p+1}} = \sum_{k=0}^{p+1}(-1)^k\alpha_{i_0 \cdots \widehat{i_k} \cdots i_{p+1}}|_{U_{i_0 \cdots i_{p+1}}}.
\]
Its cohomology is the \emph{\v{C}ech cohomology} $\check{H}^p(\mathcal{U}, \mathcal{F})$. Choosing a well-ordering of $I$, $\check{C}^p(\mathcal{U}, \mathcal{F}) \cong \prod_{i_0 < \cdots < i_p}\mathcal{F}(U_{i_0 \cdots i_p})$.
For a sheaf, $\check{H}^0(\mathcal{U}, \mathcal{F}) = \mathcal{F}(X)$ by the sheaf axioms. The \emph{sheafified \v{C}ech complex} $\check{\mathcal{C}}^p(\mathcal{U}, \mathcal{F})$ is the sheaf
$V \mapsto \check{C}^p(\mathcal{U} \cap V, \mathcal{F}|_V)$, with augmentation $\mathcal{F} \to \check{\mathcal{C}}^0$ given by restriction.
\end{definition}

\begin{lemma}\label{sheaves:lemma-cech-resolution}
For every sheaf $\mathcal{F}$ and open cover $\mathcal{U}$, the sequence $0 \to \mathcal{F} \to \check{\mathcal{C}}^0(\mathcal{U}, \mathcal{F}) \to \check{\mathcal{C}}^1(\mathcal{U}, \mathcal{F}) \to \cdots$ is exact.
If $\mathcal{F}$ is flasque, then $\check{H}^p(\mathcal{U}, \mathcal{F}) = 0$ for $p > 0$.
\end{lemma}

\begin{proof}
Let $x \in X$ and choose $j$ with $x \in U_j$. For open $V \subseteq U_j$ define $k \colon \check{C}^p(\mathcal{U} \cap V, \mathcal{F}) \to \check{C}^{p-1}(\mathcal{U} \cap V, \mathcal{F})$ by
$(k\alpha)_{i_0 \cdots i_{p-1}} = \alpha_{ji_0 \cdots i_{p-1}}$, which makes sense since $V \cap U_{ji_0 \cdots i_{p-1}} = V \cap U_{i_0 \cdots i_{p-1}}$, and $k = $ restriction to $V \cap U_j = V$ in
degree $0$. Then
$(dk\alpha)_{i_0 \cdots i_p} = \sum_k(-1)^k\alpha_{ji_0 \cdots \widehat{i_k} \cdots i_p}$ and $(kd\alpha)_{i_0 \cdots i_p} = \alpha_{i_0 \cdots i_p} - \sum_k(-1)^k\alpha_{ji_0 \cdots \widehat{i_k} \cdots i_p}$, so
$dk + kd = \id$, and similarly in degree $0$ with the augmentation. Passing to the colimit over $V \ni x$, the complex of stalks is contractible, hence exact.

If $\mathcal{F}$ is flasque, each $\check{\mathcal{C}}^p(\mathcal{U}, \mathcal{F})$ is a product of sheaves $(j_{i_0 \cdots i_p})_*(\mathcal{F}|_{U_{i_0 \cdots i_p}})$, hence flasque
(Lemma~\ref{sheaves:lemma-injective-flasque}); so the resolution is flasque, its global sections $\check{C}^\bullet(\mathcal{U}, \mathcal{F})$ compute $H^\bullet(X, \mathcal{F})$
(Proposition~\ref{sheaves:proposition-flasque}), and this vanishes in positive degrees.
\end{proof}

\begin{theorem}[Leray]\label{sheaves:theorem-leray}
Let $\mathcal{U}$ be an open cover of $X$ such that $H^q(U_{i_0 \cdots i_p}, \mathcal{F}) = 0$ for all $q > 0$ and all finite intersections. Then
$\check{H}^p(\mathcal{U}, \mathcal{F}) \cong H^p(X, \mathcal{F})$ for all $p$.
\end{theorem}

\begin{proof}
Let $\mathcal{F} \to I^\bullet$ be an injective resolution and consider the first-quadrant double complex $C^{p,q} = \check{C}^p(\mathcal{U}, I^q)$. For fixed $q$, the row
$\check{C}^\bullet(\mathcal{U}, I^q)$ is exact in positive degrees with kernel $I^q(X)$ in degree $0$, by Lemmas~\ref{sheaves:lemma-cech-resolution}
and~\ref{sheaves:lemma-injective-flasque}. For fixed $p$, the column is $\prod_{i_0 < \cdots < i_p}I^\bullet(U_{i_0 \cdots i_p})$; since $I^\bullet|_U$ is an injective
resolution of $\mathcal{F}|_U$ (Lemma~\ref{sheaves:lemma-adjoints-injectives}) and products of abelian groups are exact, its cohomology is
$\prod H^q(U_{i_0 \cdots i_p}, \mathcal{F})$, which is $0$ for $q > 0$ and $\check{C}^p(\mathcal{U}, \mathcal{F})$ for $q = 0$. By Lemma~\ref{homological-algebra:lemma-acyclic-assembly},
applied once to rows and once to columns, both $\Gamma(X, I^\bullet)$ and $\check{C}^\bullet(\mathcal{U}, \mathcal{F})$ are quasi-isomorphic to the total complex. Taking cohomology gives the
theorem.
\end{proof}

\begin{remark}\label{sheaves:remark-cech-limit}
Refinements of covers induce maps on \v{C}ech cohomology, well defined up to homotopy, and the colimit over all covers is the \v{C}ech cohomology $\check{H}^p(X, \mathcal{F})$. There
is always a natural map $\check{H}^p(X, \mathcal{F}) \to H^p(X, \mathcal{F})$, which is an isomorphism for $p \leq 1$, and for all $p$ if $X$ is paracompact Hausdorff
\cite[Chapter II, Section 5]{godement-faisceaux}. Cohomology classes in $\check{H}^1(\mathcal{U}, \mathcal{F})$ are represented by cocycles $(g_{ij})$ with $g_{ij} + g_{jk} = g_{ik}$; for
nonabelian sheaves of groups such as $\GL_n(\mathcal{O}_X)$ the same formula (written multiplicatively) classifies locally free sheaves of rank $n$ (Chapter~\ref{vector-bundles}).
\end{remark}

\section{Comparison with singular cohomology}\label{sheaves:section-singular}

\begin{theorem}\label{sheaves:theorem-singular-comparison}
Let $X$ be a paracompact Hausdorff, locally contractible space (for instance a manifold or a CW complex) and $A$ an abelian group. Then $H^p(X, A_X) \cong H^p(X; A)$, the singular
cohomology, naturally in $X$ (and compatibly with cup products \cite[Chapter III]{bredon-sheaf}).
\end{theorem}

\begin{proof}[Sketch of proof]
Let $\mathcal{S}^p$ be the sheafification of the presheaf $U \mapsto C^p(U; A)$ of singular cochains. Since every point has arbitrarily small contractible neighbourhoods, the stalks of
the cohomology presheaves $U \mapsto H^p(U; A)$ vanish for $p > 0$ and equal $A$ for $p = 0$, so $0 \to A_X \to \mathcal{S}^\bullet$ is a resolution (sheafification preserves stalks and
stalks are exact). The omitted verifications are that the sheaves $\mathcal{S}^p$ are soft on a paracompact space, and that
$\Gamma(X, \mathcal{S}^\bullet) = C^\bullet(X; A)/C^\bullet_0(X; A)$, where $C^\bullet_0$ consists of cochains vanishing on sufficiently small simplices of some open cover, which is acyclic by
the subdivision argument of Proposition~\ref{singular-homology:proposition-small-chains}; see \cite[Chapter III]{bredon-sheaf}. Then
Theorem~\ref{sheaves:theorem-soft-acyclic} gives $H^p(X, A_X) \cong H^p(C^\bullet(X; A)/C_0^\bullet) \cong H^p(X; A)$.
\end{proof}

\begin{example}\label{sheaves:example-constant-sheaf-cohomology}
$H^p(\RR^n, \ZZ) = 0$ for $p > 0$, $H^p(S^n, \ZZ)$ is $\ZZ$ for $p = 0, n$, and on a closed oriented surface of genus $g$, $H^1(\Sigma_g, \ZZ) \cong \ZZ^{2g}$. On spaces that are not locally
contractible sheaf and singular cohomology can differ: for the topologist's sine curve $T$ of Counterexample~\ref{general-topology:counterexample-sine-curve}, singular cohomology sees two
path components, $H^0(T; \ZZ) = \ZZ^2$, while $H^0(T, \ZZ_T) = \ZZ$ since $T$ is connected.
\end{example}

\begin{example}\label{sheaves:example-mobius-local-system}
Cover $S^1$ by two open arcs $U$, $V$ with $U \cap V = W_1 \sqcup W_2$, two disjoint arcs. Let $\mathcal{L}$ be the local system of rank one with monodromy $-1$, obtained by gluing $\ZZ_U$ and $\ZZ_V$ by the identity over
$W_1$ and by $-1$ over $W_2$ (Proposition~\ref{sheaves:proposition-gluing} and Definition~\ref{sheaves:definition-local-system}). The arcs and their intersections are homeomorphic to intervals, so the constant sheaf $\ZZ$ has no
higher cohomology on them (Theorem~\ref{sheaves:theorem-singular-comparison}), and Leray's theorem (Theorem~\ref{sheaves:theorem-leray}) lets us compute with the cover $\{U, V\}$. In the trivialisation over $V$,
the \v{C}ech complex is $\ZZ^2 \to \ZZ^2$, $(a, b) \mapsto (b - a, b + a)$, one coordinate for each $W_i$. Its kernel is $0$ and its image is $\{(x, y) : x \equiv y \pmod 2\}$, so
\[
H^0(S^1, \mathcal{L}) = 0, \qquad H^1(S^1, \mathcal{L}) \cong \ZZ/2 .
\]
For the constant sheaf the differential is $(a, b) \mapsto (b - a, b - a)$, which gives $H^0 = H^1 = \ZZ$, in accordance with Theorem~\ref{sheaves:theorem-singular-comparison}. With coefficients in $\QQ$ the twisted local system
has $H^0 = H^1 = 0$, since the determinant $2$ of the differential becomes invertible: a nontrivial local system can have vanishing rational cohomology on a space whose ordinary cohomology does not vanish.
\end{example}

\section{Hypercohomology and the Leray spectral sequence}\label{sheaves:section-hypercohomology}

\begin{proposition}\label{sheaves:proposition-hypercohomology}
Let $\mathcal{K}^\bullet$ be a complex of sheaves with $\mathcal{K}^p = 0$ for $p < 0$.
\begin{enumerate}
\item A quasi-isomorphism $\mathcal{K}^\bullet \to \mathcal{L}^\bullet$ induces isomorphisms on hypercohomology.
\item There are convergent spectral sequences $E_1^{p,q} = H^q(X, \mathcal{K}^p) \Rightarrow \mathbb{H}^{p+q}(X, \mathcal{K}^\bullet)$ and
$E_2^{p,q} = H^p(X, \mathcal{H}^q(\mathcal{K}^\bullet)) \Rightarrow \mathbb{H}^{p+q}(X, \mathcal{K}^\bullet)$, where $\mathcal{H}^q$ are the cohomology sheaves.
\item If every $\mathcal{K}^p$ is $\Gamma$-acyclic, $\mathbb{H}^p(X, \mathcal{K}^\bullet) = H^p(\Gamma(X, \mathcal{K}^\bullet))$. If $\mathcal{F} \to \mathcal{K}^\bullet$ is a resolution, then
$\mathbb{H}^p(X, \mathcal{K}^\bullet) = H^p(X, \mathcal{F})$.
\end{enumerate}
\end{proposition}

\begin{proof}
(1) Quasi-isomorphisms are isomorphisms in the derived category, and $R\Gamma$ is a functor on it. (2) is Proposition~\ref{homological-algebra:proposition-hypercohomology}.
(3) The first statement is Lemma~\ref{homological-algebra:lemma-acyclic-complexes}, and the second follows from (1) applied to $\mathcal{F} \to \mathcal{K}^\bullet$, with $\mathcal{F}$ in degree $0$.
\end{proof}

\begin{remark}\label{sheaves:remark-holomorphic-de-rham}
On a complex manifold $X$ the holomorphic de Rham complex $\Omega^\bullet_X$ is a resolution of $\CC_X$ (the holomorphic Poincar\'e lemma), so $H^k(X; \CC) \cong \mathbb{H}^k(X, \Omega^\bullet_X)$,
and the first spectral sequence of Proposition~\ref{sheaves:proposition-hypercohomology} becomes the \emph{Hodge-to-de Rham} (Fr\"olicher) spectral sequence
$E_1^{p,q} = H^q(X, \Omega^p_X) \Rightarrow H^{p+q}(X; \CC)$. Its degeneration for compact K\"ahler manifolds is a form of the Hodge decomposition
(Chapter~\ref{kahler}).
\end{remark}

\begin{theorem}[Leray spectral sequence]\label{sheaves:theorem-leray-spectral-sequence}
Let $f \colon X \to Y$ be continuous and $\mathcal{F}$ a sheaf of abelian groups on $X$. Then $R^qf_*\mathcal{F}$ is the sheafification of the presheaf $V \mapsto H^q(f^{-1}V, \mathcal{F})$,
and there is a convergent spectral sequence $E_2^{p,q} = H^p(Y, R^qf_*\mathcal{F}) \Rightarrow H^{p+q}(X, \mathcal{F})$.
\end{theorem}

\begin{proof}
Let $\mathcal{F} \to I^\bullet$ be an injective resolution. Then $R^qf_*\mathcal{F} = \mathcal{H}^q(f_*I^\bullet)$, the sheafification of the presheaf of cohomology groups
$V \mapsto H^q(I^\bullet(f^{-1}V))$, and $H^q(I^\bullet(f^{-1}V)) = H^q(f^{-1}V, \mathcal{F})$ because $I^\bullet|_{f^{-1}V}$ is an injective resolution
(Lemma~\ref{sheaves:lemma-adjoints-injectives}). Since $\Gamma(X, -) = \Gamma(Y, -) \circ f_*$ and $f_*$ preserves injectives, which are $\Gamma(Y, -)$-acyclic, the Grothendieck
spectral sequence (Theorem~\ref{homological-algebra:theorem-grothendieck}) applies.
\end{proof}

\begin{example}\label{sheaves:example-leray}
If $f$ is a closed embedding, $f_*$ is exact (compute stalks as in Proposition~\ref{sheaves:proposition-extension-by-zero}), so $H^p(Y, f_*\mathcal{F}) \cong H^p(X, \mathcal{F})$. For a locally
trivial fibre bundle $f \colon E \to B$ with fibre $F$ over a locally contractible paracompact base, $R^qf_*\ZZ_E$ is the local system with stalks $H^q(F; \ZZ)$, and the Leray spectral sequence
becomes the \emph{Serre spectral sequence} $E_2^{p,q} = H^p(B, \mathcal{H}^q(F)) \Rightarrow H^{p+q}(E; \ZZ)$ with local coefficients \cite{bredon-sheaf}.
\end{example}

\section{Cohomology with compact supports}\label{sheaves:section-compact-supports}

\begin{definition}\label{sheaves:definition-compact-supports}
Let $X$ be locally compact Hausdorff. The \emph{support} of a section $s \in \mathcal{F}(X)$ is the closed set $\{x : s_x \neq 0\}$. Let $\Gamma_c(X, \mathcal{F})$ be the group of sections with
compact support, a left exact functor, and $H^i_c(X, \mathcal{F}) = R^i\Gamma_c(X, \mathcal{F})$.
\end{definition}

\begin{remark}\label{sheaves:remark-compact-supports}
Flasque sheaves are $\Gamma_c$-acyclic, and on locally compact Hausdorff spaces that are countable at infinity, so are soft sheaves; so $H^*_c$ can be computed by fine resolutions. For
a manifold $M$ and an abelian group $A$, $H^i_c(M, A_M)$ agrees with the compactly supported singular cohomology of Definition~\ref{poincare-duality:definition-compact-supports}, and
Poincar\'e duality generalises to \emph{Verdier duality} $R\Gamma_c(X, \mathcal{F})^\vee \cong R\Hom(\mathcal{F}, \omega_X)$ in the derived category, with $\omega_X = \mathit{or}_X[n]$ for an
$n$-manifold \cite[Chapter III]{ks-sheaves-manifolds}, \cite{iversen-sheaves}.
\end{remark}

\section{Exercises}\label{sheaves:section-exercises}

\subsection{Sheaves on Topological Spaces}\label{sheaves:subsection-exercises-sheaves-on-topological-spaces}

\begin{exercise}\label{sheaves:exercise-constant-sheaf}
Show that the sheafification of the constant presheaf with value $A$ is the constant sheaf $A_X$ of locally constant functions, and that $A_X = \pi^{-1}A$ for
$\pi \colon X \to \mathrm{pt}$.
\end{exercise}

\begin{solution}
The stalks of the constant presheaf are all $A$, and a section of the sheafification is a function $U \to A$ locally equal to a constant, that is, a locally constant function.
For the second statement, $\pi_p^{-1}A(U) = A$ for $U \neq \emptyset$.
\end{solution}

\begin{exercise}\label{sheaves:exercise-skyscraper-adjunction}
Let $i \colon \{x\} \to X$. Show that $x_*A = i_*A$ and $\Hom(\mathcal{F}, x_*A) \cong \Hom(\mathcal{F}_x, A)$ for every sheaf $\mathcal{F}$ of abelian groups.
\end{exercise}

\begin{solution}
The first statement is the definition. By Proposition~\ref{sheaves:proposition-images}(2) and (1),
$\Hom(\mathcal{F}, i_*A) \cong \Hom(i^{-1}\mathcal{F}, A) = \Hom(\mathcal{F}_x, A)$, since a sheaf on a point is a single group.
\end{solution}

\begin{exercise}\label{sheaves:exercise-direct-image-stalk}
Give an example of a continuous map $f$ and a sheaf $\mathcal{F}$ such that $(f_*\mathcal{F})_y$ is not isomorphic to the global sections of $\mathcal{F}$ on the fibre $f^{-1}(y)$.
\end{exercise}

\begin{exercise}\label{sheaves:exercise-epi-not-surjective}
Let $X = \CC \setminus \{0\}$ with the sheaf $\mathcal{O}$ of holomorphic functions and $\mathcal{O}^\times$ of nowhere vanishing ones. Accepting that holomorphic functions have local
logarithms of nonvanishing functions but $z$ has no global logarithm, show that $\exp \colon \mathcal{O} \to \mathcal{O}^\times$, $f \mapsto e^{2\pi i f}$, is an epimorphism of sheaves but not
surjective on global sections, and identify its kernel.
\end{exercise}

\begin{exercise}\label{sheaves:exercise-restriction-open-adjoint}
For $j \colon U \to X$ open, show that $j_!$ is left adjoint to $j^{-1}$ on sheaves of abelian groups, and that $j^{-1}$ is left adjoint to $j_*$.
\end{exercise}

\subsection{Sheaf Cohomology}\label{sheaves:subsection-exercises-sheaf-cohomology}

\begin{exercise}\label{sheaves:exercise-cech-circle}
Compute $\check{H}^\bullet(\mathcal{U}, \ZZ_{S^1})$ for the cover of $S^1$ by two open arcs whose intersection has two components, and check with
Theorems~\ref{sheaves:theorem-leray} and~\ref{sheaves:theorem-singular-comparison} that $H^1(S^1, \ZZ) \cong \ZZ$.
\end{exercise}

\begin{solution}
With $I = \{0, 1\}$, $\check{C}^0 = \ZZ_{S^1}(U_0) \oplus \ZZ_{S^1}(U_1) = \ZZ^2$, $\check{C}^1 = \ZZ_{S^1}(U_0 \cap U_1) = \ZZ^2$ and $d(a, b) = (b - a, b - a)$. So $\check{H}^0 = \ZZ$ and
$\check{H}^1 = \ZZ^2/\{(c, c)\} \cong \ZZ$. The arcs and the components of their intersection are contractible, hence have no higher cohomology with constant coefficients, so Leray's theorem
applies.
\end{solution}

\begin{exercise}\label{sheaves:exercise-skyscraper-flasque}
Show that skyscraper sheaves are flasque, and that $H^i(X, x_*A) = 0$ for $i > 0$.
\end{exercise}

\begin{solution}
Restriction maps are identities or maps to $0$. Flasque sheaves are acyclic (Proposition~\ref{sheaves:proposition-flasque}).
\end{solution}

\begin{exercise}\label{sheaves:exercise-noetherian-vanishing}
Let $X$ be an irreducible space (every nonempty open subset is dense). Show that every constant sheaf on $X$ is flasque, and deduce $H^i(X, A_X) = 0$ for $i > 0$. (This applies to
algebraic varieties with the Zariski topology, and shows that sheaf cohomology of constant sheaves in that topology does not see the classical topology.)
\end{exercise}

\begin{solution}
Nonempty open subsets of an irreducible space are connected (two disjoint nonempty open subsets would contradict density), so locally constant functions on them are constant, and
$A_X(U) = A$ for $U \neq \emptyset$, with identity restriction maps.
\end{solution}

\begin{exercise}\label{sheaves:exercise-mayer-vietoris-sheaves}
For open $U, V \subseteq X$ and a sheaf $\mathcal{F}$, construct a long exact Mayer--Vietoris sequence
$\cdots \to H^i(U \cup V, \mathcal{F}) \to H^i(U, \mathcal{F}) \oplus H^i(V, \mathcal{F}) \to H^i(U \cap V, \mathcal{F}) \to H^{i+1}(U \cup V, \mathcal{F}) \to \cdots$.
(Hint: for flasque $I$, $0 \to I(U \cup V) \to I(U) \oplus I(V) \to I(U \cap V) \to 0$ is exact.)
\end{exercise}
